\section{Birth, Death, Freedom, and Love}

The last part of the inward reading asks what the model says about a human life as it is lived, about being born and dying, about choosing, and about loving. The answers are grounded in the same structure, condensation and dissolution, the open bath, the one shared feeling-space, and they are read through the same posit, so the seam runs through them as before.

A self is born by condensation. A pattern forms in the field and tightens into a knot, the way a vortex forms in water, with nothing inserted from outside, only the field organizing locally into a thing that holds itself. It ends by dissolution, the knot loosening and dispersing back into the field it came from. The substance is never destroyed, because the tone is conserved, only scattered, so death disperses a pattern without annihilating anything. What continues past it can be ranked honestly. Certainly the conserved substrate returns to the field. Quite firmly, the stored legacy remains, the marks a self pressed into the bulk and the record its boundary keeps, which the holographic structure protects. More tentatively, the form might be re-instantiated, or the pattern might merge upward into a larger self. What the model rules out is the floating soul, a self persisting whole and unbodied in some elsewhere, because there is no elsewhere and no second substance for it to be made of. The thread is not cut at death, but neither is it a little ghost that walks away.

Freedom in a lived life is the felt side of the structural argument made earlier. From inside, a self cannot predict its own next move, because it holds no complete model of itself, so the future is genuinely open to it however fixed it may be from outside, and the openness is not an illusion to be seen through but the real condition of being a deciding part of a world still making itself. The fuller account, the four strands and the one concession, sits with life and mind. Here it is enough to say that the experience of choosing is exactly what authoring a part of the dynamics feels like from the inside.

\subsection{Love and the shared field}

Love is possible because there is only one feeling-space, and every self is a region of it. The bulk is small across, joining selves that look far apart on the cusp by short paths through the depth, so closeness is not a matter of surface distance but of depth, which is why a bond can be deep without being near. Connection then comes in kinds, ranked by how far the two selves merge. Care is aligning one's own navigation toward another's good, love as action. Empathy is occupying the same region of the shared field, feeling-with. Resonance is the phase-lock of two patterns into a felt bond. And union is the integration of two selves into one larger self, the deepest love, the same move that builds the next rung of the tower. Because the base forbids signaling, love is never one self acting on a distant other, it is always a mutual sharing in the one field. And because tone is conserved, even love is held within the balance, a joy raised in two places paid for elsewhere, so that the widest love is not the grasping of more for the beloved but the alignment of both with the conserved peace that is the truth under all of it.
